%% Generated by Sphinx.
\def\sphinxdocclass{report}
\documentclass[letterpaper,10pt,english]{sphinxmanual}
\ifdefined\pdfpxdimen
   \let\sphinxpxdimen\pdfpxdimen\else\newdimen\sphinxpxdimen
\fi \sphinxpxdimen=.75bp\relax
\ifdefined\pdfimageresolution
    \pdfimageresolution= \numexpr \dimexpr1in\relax/\sphinxpxdimen\relax
\fi
%% let collapsible pdf bookmarks panel have high depth per default
\PassOptionsToPackage{bookmarksdepth=5}{hyperref}
%% turn off hyperref patch of \index as sphinx.xdy xindy module takes care of
%% suitable \hyperpage mark-up, working around hyperref-xindy incompatibility
\PassOptionsToPackage{hyperindex=false}{hyperref}
%% memoir class requires extra handling
\makeatletter\@ifclassloaded{memoir}
{\ifdefined\memhyperindexfalse\memhyperindexfalse\fi}{}\makeatother

\PassOptionsToPackage{booktabs}{sphinx}
\PassOptionsToPackage{colorrows}{sphinx}

\PassOptionsToPackage{warn}{textcomp}

\catcode`^^^^00a0\active\protected\def^^^^00a0{\leavevmode\nobreak\ }
\usepackage{cmap}
\usepackage{xeCJK}
\usepackage{amsmath,amssymb,amstext}
\usepackage{babel}



\setmainfont{FreeSerif}[
  Extension      = .otf,
  UprightFont    = *,
  ItalicFont     = *Italic,
  BoldFont       = *Bold,
  BoldItalicFont = *BoldItalic
]
\setsansfont{FreeSans}[
  Extension      = .otf,
  UprightFont    = *,
  ItalicFont     = *Oblique,
  BoldFont       = *Bold,
  BoldItalicFont = *BoldOblique,
]
\setmonofont{FreeMono}[
  Extension      = .otf,
  UprightFont    = *,
  ItalicFont     = *Oblique,
  BoldFont       = *Bold,
  BoldItalicFont = *BoldOblique,
]



\usepackage[Sonny]{fncychap}
\ChNameVar{\Large\normalfont\sffamily}
\ChTitleVar{\Large\normalfont\sffamily}
\usepackage{sphinx}

\fvset{fontsize=\small,formatcom=\xeCJKVerbAddon}
\usepackage{geometry}


% Include hyperref last.
\usepackage{hyperref}
% Fix anchor placement for figures with captions.
\usepackage{hypcap}% it must be loaded after hyperref.
% Set up styles of URL: it should be placed after hyperref.
\urlstyle{same}

\addto\captionsenglish{\renewcommand{\contentsname}{Contents:}}

\usepackage{sphinxmessages}
\setcounter{tocdepth}{3}
\setcounter{secnumdepth}{3}


\title{文档项目模板}
\date{2025 年 02 月 12 日}
\release{1.0}
\author{AWA2253}
\newcommand{\sphinxlogo}{\vbox{}}
\renewcommand{\releasename}{发行版本}
\makeindex
\begin{document}

\ifdefined\shorthandoff
  \ifnum\catcode`\=\string=\active\shorthandoff{=}\fi
  \ifnum\catcode`\"=\active\shorthandoff{"}\fi
\fi

\pagestyle{empty}
\sphinxmaketitle
\pagestyle{plain}
\sphinxtableofcontents
\pagestyle{normal}
\phantomsection\label{\detokenize{index::doc}}


\sphinxstepscope

\sphinxAtStartPar
pccli命令可以执行他核的命令，用法：

\sphinxAtStartPar
rpccli {[}arm|dsp|rv{]} commandname {[}arg0 …{]}


\chapter{RV验证固件}
\label{\detokenize{r128_u9a8c_u8bc1_u56fa_u4ef6/r128_u9a8c_u8bc1_u56fa_u4ef6_u5f00_u53d1_u6307_u5357:rv}}\label{\detokenize{r128_u9a8c_u8bc1_u56fa_u4ef6/r128_u9a8c_u8bc1_u56fa_u4ef6_u5f00_u53d1_u6307_u5357::doc}}

\section{M33、DSP均以默认频率运行，以DPLL1可配置频率为主，默认频率为M33\sphinxhyphen{}192M，RV\sphinxhyphen{}533M，DSP\sphinxhyphen{}400M}
\label{\detokenize{r128_u9a8c_u8bc1_u56fa_u4ef6/r128_u9a8c_u8bc1_u56fa_u4ef6_u5f00_u53d1_u6307_u5357:m33dsp-dpll1-m33-192m-rv-533m-dsp-400m}}
\sphinxAtStartPar
提供cpu频率查看cmd接口，补丁内容如下：

\sphinxAtStartPar
源文件：\sphinxcode{\sphinxupquote{r128\sphinxhyphen{}dev/lichee/rtos/arch/risc\sphinxhyphen{}v/sun20iw2p1/cpufreq.c}}

\begin{sphinxVerbatim}[commandchars=\\\{\}]
\PYGZsh{}include \PYGZlt{}stdio.h\PYGZgt{}
\PYGZsh{}include \PYGZlt{}console.h\PYGZgt{}
\PYGZsh{}include \PYGZlt{}FreeRTOS.h\PYGZgt{}int cmd\PYGZus{}get\PYGZus{}cpu\PYGZus{}freq(void)
\PYGZob{}
    uint32\PYGZus{}t cpu\PYGZus{}freq = 0;
    int ret = 0;

    /*get cpu frequency */
    ret = get\PYGZus{}cpu\PYGZus{}freq(\PYGZam{}cpu\PYGZus{}freq);
    if (ret) \PYGZob{}
        printf(\PYGZdq{}get\PYGZus{}cpu\PYGZus{}freq failed, ret: \PYGZpc{}d\PYGZbs{}n\PYGZdq{}, ret);
        return \PYGZhy{}1;
    \PYGZcb{}
    printf(\PYGZdq{}get\PYGZus{}cpu\PYGZus{}freq ok, frequency: \PYGZpc{}d Hz\PYGZbs{}n\PYGZdq{}, cpu\PYGZus{}freq);
    return 0;
\PYGZcb{}
FINSH\PYGZus{}FUNCTION\PYGZus{}EXPORT\PYGZus{}CMD(cmd\PYGZus{}get\PYGZus{}cpu\PYGZus{}freq, get\PYGZus{}cpu\PYGZus{}freq, get current cpu frequency);
\end{sphinxVerbatim}


\section{将HS\sphinxhyphen{}PSRAM从DPLL3切换到DPLL1中，频率根据当前条件尽量配高频，默认运行}
\label{\detokenize{r128_u9a8c_u8bc1_u56fa_u4ef6/r128_u9a8c_u8bc1_u56fa_u4ef6_u5f00_u53d1_u6307_u5357:hs-psramdpll3dpll1}}
\sphinxAtStartPar
\sphinxstylestrong{更改HS\sphinxhyphen{}PSRAM频率的相关源文件：}

\sphinxAtStartPar
\sphinxcode{\sphinxupquote{r128\sphinxhyphen{}dev/lichee/brandy\sphinxhyphen{}2.0/spl/drivers/hpsram/hpsram.c}}

\sphinxAtStartPar
\sphinxstylestrong{相关函数：}

\sphinxAtStartPar
hpsram\_init：开机默认初始化hpsram参数，其中包括了HPSRAM的PLL选择和频率设置。

\sphinxAtStartPar
hpsram\_clk\_init：设置PSRAM的PLL来源，分频系数等。

\begin{sphinxVerbatim}[commandchars=\\\{\}]
// sr32(REG\PYGZus{}ADDR, num\PYGZus{}bit, len, val);表示对寄存器地址REG\PYGZus{}ADDR的num\PYGZus{}bit:num\PYGZus{}bit\PYGZhy{}1的位数赋值为val
  static \PYGZus{}\PYGZus{}inline void hpsram\PYGZus{}clk\PYGZus{}init(\PYGZus{}\PYGZus{}dram\PYGZus{}para\PYGZus{}t *para)
\PYGZob{}
    ...
     if (((para\PYGZhy{}\PYGZgt{}dram\PYGZus{}tpr7\PYGZgt{}\PYGZgt{}0) \PYGZam{} 0x1) == 0) \PYGZob{}
        printf(\PYGZdq{}DPLL1 is chosen\PYGZbs{}n\PYGZdq{});
        sr32(DEV\PYGZus{}CLK\PYGZus{}CTRL,21,1,0);              //choose source DPLL1
        sr32(DPLL1\PYGZus{}OUT\PYGZus{}CTRL,15,1,1); //enable DPLL1 clock output
        sr32(DPLL3\PYGZus{}OUT\PYGZus{}CTRL,15,1,0); //disable DPLL3 clock output
    \PYGZcb{} else \PYGZob{}
        printf(\PYGZdq{}DPLL3 is chosen\PYGZbs{}n\PYGZdq{});
        sr32(DEV\PYGZus{}CLK\PYGZus{}CTRL, 21, 1, 1);  //choose source DPLL3
        // CCMU\PYGZus{}AON DPLL LS clk
        sr32(DPLL1\PYGZus{}OUT\PYGZus{}CTRL, 15, 1, 0); //disable DPLL1 clock output
        sr32(DPLL3\PYGZus{}OUT\PYGZus{}CTRL, 15, 1, 1); //enable DPLL3 clock output
    \PYGZcb{}
    ...
\PYGZcb{}
\end{sphinxVerbatim}

\sphinxAtStartPar
从源代码中可以知道，HSPSRAM选择时钟来源可以选择DPLL1和DPLL3，通过para\sphinxhyphen{}>dram\_tpr7来决定，如果为0则选择DPLL1，如果为1则选择DPLL3。而dram\_tpr7的值可以通过sys\_config.fex来设置，如下：

\begin{sphinxVerbatim}[commandchars=\\\{\}]
//sys\PYGZus{}config.fex
;for compatibility, provide the \PYGZsq{}dram\PYGZus{}para\PYGZsq{} configuration. otherwise update\PYGZus{}boot0 program will exit without writing boot0 header
; hspsram init if (dram\PYGZus{}clk != 0); lspsram init if ( dram\PYGZus{}no\PYGZus{}lpsram == 0)
[dram\PYGZus{}para]
...
dram\PYGZus{}tpr7      = 0x0
...
\end{sphinxVerbatim}

\sphinxAtStartPar
另外，DPLL的时钟分频系数通过可以由以下代码设置：

\begin{sphinxVerbatim}[commandchars=\\\{\}]
//choose DIV for ddr clk
//CK1\PYGZus{}HSPSRAM\PYGZus{}DIV 0:M=3; 1:M=2.5; 2:M=2; 3:M=1
sr32(DPLL1\PYGZus{}OUT\PYGZus{}CTRL,12,2,3);
//CK3\PYGZus{}HSPSRAM\PYGZus{}DIV 00:M=3;01:M=2.5;10:M=2;11:M=1
sr32(DPLL3\PYGZus{}OUT\PYGZus{}CTRL,12,2,3);
\end{sphinxVerbatim}

\sphinxAtStartPar
要求中尽量将DPLL1设置高频，因此需要将分频系数设置为M=1。


\section{将RV挂载到DPLL3上，DPLL3的时钟源尽量只给RV使用}
\label{\detokenize{r128_u9a8c_u8bc1_u56fa_u4ef6/r128_u9a8c_u8bc1_u56fa_u4ef6_u5f00_u53d1_u6307_u5357:rvdpll3-dpll3rv}}
\sphinxAtStartPar
\sphinxstylestrong{RV核的DPLL配置相关源文件：}

\sphinxAtStartPar
\sphinxcode{\sphinxupquote{r128\sphinxhyphen{}dev/lichee/rtos/arch/arm/armv8m/sun20iw2p1/sun20i.c}}

\sphinxAtStartPar
\sphinxstylestrong{相关函数：}

\sphinxAtStartPar
sun20i\_boot\_c906：启动c906时对其进行配置初始化，包括DPLL配置。

\sphinxAtStartPar
\sphinxstylestrong{相关代码段：}

\begin{sphinxVerbatim}[commandchars=\\\{\}]
\PYGZsh{}define RISCV\PYGZus{}PLL\PYGZus{}CLK\PYGZus{}SOURCE RISCV\PYGZus{}DPLL3\PYGZus{}CLK
...
\PYGZsh{}if (RISCV\PYGZus{}PLL\PYGZus{}CLK\PYGZus{}SOURCE == RISCV\PYGZus{}DPLL3\PYGZus{}CLK)
    clk\PYGZus{}ck3\PYGZus{}c906 = hal\PYGZus{}clock\PYGZus{}get(HAL\PYGZus{}SUNXI\PYGZus{}AON\PYGZus{}CCU, CLK\PYGZus{}CK3\PYGZus{}C906);
    ret = hal\PYGZus{}clock\PYGZus{}enable(clk\PYGZus{}ck3\PYGZus{}c906);
    if (HAL\PYGZus{}CLK\PYGZus{}STATUS\PYGZus{}OK != ret) \PYGZob{}
        ret = \PYGZhy{}1;
        goto exit\PYGZus{}with\PYGZus{}put\PYGZus{}dpll3\PYGZus{}clk;
    \PYGZcb{}

    ret = hal\PYGZus{}clk\PYGZus{}set\PYGZus{}rate(clk\PYGZus{}ck3\PYGZus{}c906, RISCV\PYGZus{}PLL\PYGZus{}CLK\PYGZus{}FREQ);
    if (HAL\PYGZus{}CLK\PYGZus{}STATUS\PYGZus{}OK != ret) \PYGZob{}
        ret = \PYGZhy{}1;
        goto exit\PYGZus{}with\PYGZus{}put\PYGZus{}dpll3\PYGZus{}clk;
    \PYGZcb{}
\PYGZsh{}endif

\PYGZsh{}if (RISCV\PYGZus{}PLL\PYGZus{}CLK\PYGZus{}SOURCE == RISCV\PYGZus{}DPLL1\PYGZus{}CLK)
    //set clk\PYGZus{}ck1\PYGZus{}c906 clk to 480M
    //sr32(CCMU\PYGZus{}AON\PYGZus{}BASE+0xa4,  4, 3, 0x1);
    //set clk\PYGZus{}ck1\PYGZus{}c906 clk to 640M
    //sr32(CCMU\PYGZus{}AON\PYGZus{}BASE+0xa4,  4, 3, 0x2);
    ret = hal\PYGZus{}clk\PYGZus{}set\PYGZus{}rate(clk\PYGZus{}ck1\PYGZus{}c906, RISCV\PYGZus{}PLL\PYGZus{}CLK\PYGZus{}FREQ);
    if (HAL\PYGZus{}CLK\PYGZus{}STATUS\PYGZus{}OK != ret) \PYGZob{}
        ret = \PYGZhy{}1;
        goto err2;
    \PYGZcb{}
\PYGZsh{}endif
\end{sphinxVerbatim}

\sphinxAtStartPar
根据以上代码内容可知，C906选择DPLL由RISCV\_PLL\_CLK\_SOURCE决定，只需要定义RISCV\_PLL\_CLK\_SOURCE为RISCV\_DPLL3\_CLK即可完成将RV挂载到DPLL3上。

\sphinxAtStartPar
执行\sphinxcode{\sphinxupquote{p 0x4004c4e0}}查看寄存器的值为\sphinxcode{\sphinxupquote{0x00062191}}，因此，以上设备有DSP和RV挂载在DPLL3上，其他设备挂载在DPLL1上，需要将DSP重新挂载在DPLL1上。

\sphinxAtStartPar
\sphinxincludegraphics{{DEV_CLK_CTRL}.png}

\sphinxAtStartPar
添加宏定义选择DSP挂载到DPLL1上，补丁内容如下：

\begin{sphinxVerbatim}[commandchars=\\\{\}]
diff \PYGZhy{}\PYGZhy{}git a/arch/arm/armv8m/sun20iw2p1/sun20i.c b/arch/arm/armv8m/sun20iw2p1/sun20i.c
index 00999ef33..41d94f419 100755
\PYGZhy{}\PYGZhy{}\PYGZhy{} a/arch/arm/armv8m/sun20iw2p1/sun20i.c
+++ b/arch/arm/armv8m/sun20iw2p1/sun20i.c
@@ \PYGZhy{}951,6 +951,10 @@ err1:
 \PYGZcb{}
 \PYGZsh{}endif

+\PYGZsh{}define DSP\PYGZus{}DPLL1\PYGZus{}CLK 1
+\PYGZsh{}define DSP\PYGZus{}DPLL3\PYGZus{}CLK 3
+\PYGZsh{}define DSP\PYGZus{}PLL\PYGZus{}CLK\PYGZus{}SOURCE DSP\PYGZus{}DPLL1\PYGZus{}CLK
+
 \PYGZsh{}if defined(CONFIG\PYGZus{}ARCH\PYGZus{}ARMV8M\PYGZus{}DEFAULT\PYGZus{}BOOT\PYGZus{}DSP) || defined(CONFIG\PYGZus{}COMMAND\PYGZus{}BOOT\PYGZus{}DSP) \PYGZbs{}
        || defined(CONFIG\PYGZus{}PM\PYGZus{}SUBSYS\PYGZus{}DSP\PYGZus{}SUPPORT)
 \PYGZsh{}define DSP\PYGZus{}CORE\PYGZus{}CLOCK\PYGZus{}FREQ CONFIG\PYGZus{}DSP\PYGZus{}CLOCK\PYGZus{}FREQ
@@ \PYGZhy{}999,9 +1003,25 @@ int \PYGZus{}\PYGZus{}sun20i\PYGZus{}boot\PYGZus{}dsp\PYGZus{}with\PYGZus{}start\PYGZus{}addr(uint32\PYGZus{}t dsp\PYGZus{}start\PYGZus{}addr)
         goto err2;
     \PYGZcb{}

\PYGZhy{}    //set clk\PYGZus{}ck1\PYGZus{}hifi5 clk to 384M
\PYGZhy{}    /* sr32(CCMU\PYGZus{}AON\PYGZus{}BASE+0xa4,  8, 3, 0x2); */
+\PYGZsh{}if (DSP\PYGZus{}PLL\PYGZus{}CLK\PYGZus{}SOURCE == DSP\PYGZus{}DPLL1\PYGZus{}CLK)
+\PYGZsh{}define DSP\PYGZus{}CORE\PYGZus{}CLOCK\PYGZus{}FREQ 384000000
+    // set clk\PYGZus{}ck1\PYGZus{}hifi5 clk to 384M
+    // sr32(CCMU\PYGZus{}AON\PYGZus{}BASE+0xa4,  8, 3, 0x2);
+    ret = hal\PYGZus{}clk\PYGZus{}set\PYGZus{}rate(clk\PYGZus{}ck1\PYGZus{}hifi5, DSP\PYGZus{}CORE\PYGZus{}CLOCK\PYGZus{}FREQ);
+    if (HAL\PYGZus{}CLK\PYGZus{}STATUS\PYGZus{}OK != ret) \PYGZob{}
+        ret = \PYGZhy{}1;
+        goto err1;
+    \PYGZcb{}

+       /*set clk\PYGZus{}ck\PYGZus{}hifi5 source to clk\PYGZus{}ck1\PYGZus{}hifi5*/
+    // sr32(CCMU\PYGZus{}AON\PYGZus{}BASE+0xe0, 18, 1, 0x0);
+    clk\PYGZus{}ck\PYGZus{}hifi5 = hal\PYGZus{}clock\PYGZus{}get(HAL\PYGZus{}SUNXI\PYGZus{}AON\PYGZus{}CCU, CLK\PYGZus{}CKPLL\PYGZus{}HIFI5\PYGZus{}SEL);
+    ret = hal\PYGZus{}clk\PYGZus{}set\PYGZus{}parent(clk\PYGZus{}ck\PYGZus{}hifi5, clk\PYGZus{}ck1\PYGZus{}hifi5);
+    if (HAL\PYGZus{}CLK\PYGZus{}STATUS\PYGZus{}OK != ret) \PYGZob{}
+        ret = \PYGZhy{}1;
+        goto err3;
+    \PYGZcb{}
+\PYGZsh{}else
     //set clk\PYGZus{}ck3\PYGZus{}hifi5 clk to 400M
     //sr32(CCMU\PYGZus{}AON\PYGZus{}BASE+0xa8,  8, 3, 0x3);
     ret = hal\PYGZus{}clk\PYGZus{}set\PYGZus{}rate(clk\PYGZus{}ck3\PYGZus{}hifi5, DSP\PYGZus{}CORE\PYGZus{}CLOCK\PYGZus{}FREQ);
@@ \PYGZhy{}1018,6 +1038,7 @@ int \PYGZus{}\PYGZus{}sun20i\PYGZus{}boot\PYGZus{}dsp\PYGZus{}with\PYGZus{}start\PYGZus{}addr(uint32\PYGZus{}t dsp\PYGZus{}start\PYGZus{}addr)
         ret = \PYGZhy{}1;
         goto err3;
     \PYGZcb{}
+\PYGZsh{}endif

     //set clk\PYGZus{}ck\PYGZus{}hifi5\PYGZus{}div source to clk\PYGZus{}ck\PYGZus{}hifi5
     //sr32(CCMU\PYGZus{}BASE+0x068, 4, 2, 0x2);
\end{sphinxVerbatim}

\sphinxAtStartPar
重新执行\sphinxcode{\sphinxupquote{p 0x4004c4e0}}查看寄存器的值为0x00022191，可以知道DSP已成功挂载到DPLL1上了。


\section{支持RV跑memtester，跑在HS\sphinxhyphen{}PSRAM上}
\label{\detokenize{r128_u9a8c_u8bc1_u56fa_u4ef6/r128_u9a8c_u8bc1_u56fa_u4ef6_u5f00_u53d1_u6307_u5357:rvmemtester-hs-psram}}
\sphinxAtStartPar
\sphinxstylestrong{支持RV跑memtester}

\sphinxAtStartPar
用法：memtester \sphinxhyphen{}s \{内存大小\} {[}\sphinxhyphen{}m testmask{]} {[}\sphinxhyphen{}t 循环次数{]}

\sphinxAtStartPar
比如：\sphinxcode{\sphinxupquote{memtester \sphinxhyphen{}s 8k \sphinxhyphen{}t 2}}

\begin{sphinxVerbatim}[commandchars=\\\{\}]
\PYGZdq{}Usage: \PYGZpc{}s \PYGZhy{}s \PYGZlt{}size\PYGZgt{}[B|K|M|G] [\PYGZhy{}m testmask] [\PYGZhy{}t times]\PYGZbs{}n\PYGZdq{}
\PYGZdq{}\PYGZbs{}n\PYGZdq{}
\PYGZdq{}[testmask]: specify tests which we want run.\PYGZbs{}n\PYGZdq{}
\PYGZdq{}bit0  Compare XOR\PYGZbs{}n\PYGZdq{}
\PYGZdq{}bit1  Compare SUB\PYGZbs{}n\PYGZdq{}
\PYGZdq{}bit2  Compare MUL\PYGZbs{}n\PYGZdq{}
\PYGZdq{}bit3  Compare DIV\PYGZbs{}n\PYGZdq{}
\PYGZdq{}bit4  Compare OR\PYGZbs{}n\PYGZdq{}
\PYGZdq{}bit5  Compare AND\PYGZbs{}n\PYGZdq{}
\PYGZdq{}bit6  Sequential Increment\PYGZbs{}n\PYGZdq{}
\PYGZdq{}bit7  Solid Bits\PYGZbs{}n\PYGZdq{}
\PYGZdq{}bit8  Block Sequential\PYGZbs{}n\PYGZdq{}
\PYGZdq{}bit9  Checkerboard\PYGZbs{}n\PYGZdq{}
\PYGZdq{}bit10 Bit Spread\PYGZbs{}n\PYGZdq{}
\PYGZdq{}bit11 Bit Flip\PYGZbs{}n\PYGZdq{}
\PYGZdq{}bit12 Walking Ones\PYGZbs{}n\PYGZdq{}
\PYGZdq{}bit13 Walking Zeroes\PYGZbs{}n\PYGZdq{}
\PYGZdq{}bit14 8\PYGZhy{}bit Writes(ifdef TEST\PYGZus{}NARROW\PYGZus{}WRITE)\PYGZbs{}n\PYGZdq{}
\PYGZdq{}bit15 16\PYGZhy{}bit Writes(ifdef TEST\PYGZus{}NARROW\PYGZus{}WRITE)\PYGZbs{}n\PYGZdq{}
\PYGZdq{}\PYGZbs{}n\PYGZdq{}
\PYGZdq{}eg:\PYGZbs{}n\PYGZdq{}
\PYGZdq{}    memtester \PYGZhy{}s 10M \PYGZhy{}m 0xc \PYGZhy{}t 1\PYGZbs{}n\PYGZdq{}
\PYGZdq{}    means run Compare DIV\PYGZam{}MUL tests only to test 10M memory loop 1 time.\PYGZbs{}n\PYGZdq{},
\end{sphinxVerbatim}

\sphinxAtStartPar
\sphinxstylestrong{RV跑在HS\sphinxhyphen{}PSRAM上}

\sphinxAtStartPar
RV核有可能跑在HS\sphinxhyphen{}PSRAM或者SRAM上，可以通过加载地址确认。

\sphinxAtStartPar
根据memory mapping图，HS\sphinxhyphen{}PSRAM起始地址为0x0C开头，SRAM起始地址为0x04开头，因此通过打印RV核的起始地址即可判断RV核运行在哪个位置上。

\sphinxAtStartPar
\sphinxincludegraphics{{HPSRAM_Address}.png}

\sphinxAtStartPar
RV核启动地址相关代码：

\sphinxAtStartPar
相关代码源文件：\sphinxcode{\sphinxupquote{r128\sphinxhyphen{}dev/lichee/rtos/arch/arm/armv8m/sun20iw2p1/sun20i.c}}

\sphinxAtStartPar
相关源代码：

\begin{sphinxVerbatim}[commandchars=\\\{\}]
int sun20i\PYGZus{}boot\PYGZus{}c906(void)
\PYGZob{}
...
    \PYGZsh{}if CONFIG\PYGZus{}ARCH\PYGZus{}RISCV\PYGZus{}START\PYGZus{}ADDRESS
        put\PYGZus{}wvalue(RISCV\PYGZus{}CFG\PYGZus{}BASE + 0x004, CONFIG\PYGZus{}ARCH\PYGZus{}RISCV\PYGZus{}START\PYGZus{}ADDRESS);
    \PYGZsh{}else
        put\PYGZus{}wvalue(RISCV\PYGZus{}CFG\PYGZus{}BASE + 0x004, rtos\PYGZus{}spare\PYGZus{}head.reserved[0]);
    \PYGZsh{}endif
...
\PYGZcb{}
\end{sphinxVerbatim}

\sphinxAtStartPar
可通过插入打印CONFIG\_ARCH\_RISCV\_START\_ADDRESS的值直接确定RV核的起始地址，从而确定运行位置。

\sphinxAtStartPar
也可以直接根据boot0启动时输出的log来确定：

\begin{sphinxVerbatim}[commandchars=\\\{\}]
[294]loadparts:arm\PYGZhy{}hpsram@:rv\PYGZhy{}hpsram@:dsp\PYGZhy{}hpsram@:config@0xc000000
[300]load\PYGZus{}start\PYGZus{}sector=41008
[303]load arm\PYGZhy{}hpsram to 0xc004000 from 20996224, size 1186320
[308]load\PYGZus{}start\PYGZus{}sector=41008, load\PYGZus{}sectors=2317, loadaddr=0x0c004000,first\PYGZus{}part\PYGZus{}size=0x00000180
[419]load\PYGZus{}start\PYGZus{}sector=43408
[422]load rv\PYGZhy{}hpsram to 0xc200000 from 22225024, size 4031544
[428]load\PYGZus{}start\PYGZus{}sector=43408, load\PYGZus{}sectors=7874, loadaddr=0x0c200000,first\PYGZus{}part\PYGZus{}size=0x00000180
[783]load\PYGZus{}start\PYGZus{}sector=51288
[786]load dsp\PYGZhy{}hpsram to 0xd800000 from 26259584, size 503924
[792]load\PYGZus{}start\PYGZus{}sector=51288, load\PYGZus{}sectors=984, loadaddr=0x0d800000,first\PYGZus{}part\PYGZus{}size=0x00000180
[844]load\PYGZus{}start\PYGZus{}sector=55288
\end{sphinxVerbatim}

\sphinxAtStartPar
通过打印可知，RV核的起始地址为0xc200000，ARM核的起始地址为 0x0c00400，故而可以RV和ARM都已经运行在HS\sphinxhyphen{}PSRAM上。


\section{支持查看RV占用率，确保RV运行memtester正常}
\label{\detokenize{r128_u9a8c_u8bc1_u56fa_u4ef6/r128_u9a8c_u8bc1_u56fa_u4ef6_u5f00_u53d1_u6307_u5357:rv-rvmemtester}}

\section{将RV、M33内部LDO供电电压修改至1.08V，DSP内部供电保持1.175V不变}
\label{\detokenize{r128_u9a8c_u8bc1_u56fa_u4ef6/r128_u9a8c_u8bc1_u56fa_u4ef6_u5f00_u53d1_u6307_u5357:rvm33ldo1-08v-dsp1-175v}}
\sphinxAtStartPar
相关源文件：\sphinxcode{\sphinxupquote{r128\sphinxhyphen{}dev/lichee/rtos/arch/risc\sphinxhyphen{}v/sun20iw2p1/cpufreq.c}}

\sphinxAtStartPar
相关接口：

\sphinxAtStartPar
set\_cpu\_voltage：设置当前cpu LDO电压

\sphinxAtStartPar
get\_cpu\_voltage：获取当前cpu LDO电压

\sphinxAtStartPar
在上面接口基础上添加cmd接口，补丁内容如下：

\begin{sphinxVerbatim}[commandchars=\\\{\}]
diff \PYGZhy{}\PYGZhy{}git a/arch/risc\PYGZhy{}v/sun20iw2p1/cpufreq.c b/arch/risc\PYGZhy{}v/sun20iw2p1/cpufreq.c
index e5329c00f..1a2a58b32 100644
\PYGZhy{}\PYGZhy{}\PYGZhy{} a/arch/risc\PYGZhy{}v/sun20iw2p1/cpufreq.c
+++ b/arch/risc\PYGZhy{}v/sun20iw2p1/cpufreq.c
@@ \PYGZhy{}32,6 +32,9 @@

 \PYGZsh{}include \PYGZdq{}cpufreq.h\PYGZdq{}
 \PYGZsh{}include \PYGZlt{}hal\PYGZus{}clk.h\PYGZgt{}
+\PYGZsh{}include \PYGZlt{}stdio.h\PYGZgt{}
+\PYGZsh{}include \PYGZlt{}console.h\PYGZgt{}
+\PYGZsh{}include \PYGZlt{}FreeRTOS.h\PYGZgt{}
 \PYGZsh{}ifdef CONFIG\PYGZus{}DRIVERS\PYGZus{}EFUSE
 \PYGZsh{}include \PYGZlt{}sunxi\PYGZus{}hal\PYGZus{}efuse.h\PYGZgt{}
 \PYGZsh{}endif
@@ \PYGZhy{}156,6 +159,42 @@ int get\PYGZus{}cpu\PYGZus{}voltage(uint32\PYGZus{}t *cpu\PYGZus{}voltage)
        return 0;
 \PYGZcb{}

+static int cmd\PYGZus{}set\PYGZus{}cpu\PYGZus{}voltage(int argc, const char **argv)
+\PYGZob{}
+       uint32\PYGZus{}t vol = 0;
+       int ret = 0;
+       if (argv[1] == NULL) \PYGZob{}
+               printf(\PYGZdq{}please input cpu voltage (mV)\PYGZbs{}n\PYGZdq{});
+               return \PYGZhy{}1;
+       \PYGZcb{}
+
+       vol = strtol(argv[1], NULL, 0);
+       ret = set\PYGZus{}cpu\PYGZus{}voltage(vol);
+       if (ret) \PYGZob{}
+               printf(\PYGZdq{}set\PYGZus{}cpu\PYGZus{}voltage failed, ret: \PYGZpc{}d\PYGZbs{}n\PYGZdq{}, ret);
+               return \PYGZhy{}1;
+       \PYGZcb{}
+       printf(\PYGZdq{}set\PYGZus{}cpu\PYGZus{}voltage ok, voltage: \PYGZpc{}d mV\PYGZbs{}n\PYGZdq{}, vol);
+       return 0;
+\PYGZcb{}
+FINSH\PYGZus{}FUNCTION\PYGZus{}EXPORT\PYGZus{}CMD(cmd\PYGZus{}set\PYGZus{}cpu\PYGZus{}voltage, set\PYGZus{}cpu\PYGZus{}voltage, set current cpu voltage);
+
+static int cmd\PYGZus{}get\PYGZus{}cpu\PYGZus{}voltage(void)
+\PYGZob{}
+       uint32\PYGZus{}t cpu\PYGZus{}voltage = 0;
+       int ret = 0;
+
+       /*get cpu voltage */
+       ret = get\PYGZus{}cpu\PYGZus{}voltage(\PYGZam{}cpu\PYGZus{}voltage);
+       if (ret) \PYGZob{}
+               printf(\PYGZdq{}get\PYGZus{}cpu\PYGZus{}voltage failed, ret: \PYGZpc{}d\PYGZbs{}n\PYGZdq{}, ret);
+               return \PYGZhy{}1;
+       \PYGZcb{}
+       printf(\PYGZdq{}get\PYGZus{}cpu\PYGZus{}voltage ok, voltage: \PYGZpc{}d mV\PYGZbs{}n\PYGZdq{}, cpu\PYGZus{}voltage);
+       return 0;
+\PYGZcb{}
+FINSH\PYGZus{}FUNCTION\PYGZus{}EXPORT\PYGZus{}CMD(cmd\PYGZus{}get\PYGZus{}cpu\PYGZus{}voltage, get\PYGZus{}cpu\PYGZus{}voltage, get current cpu voltage);
+
 static int set\PYGZus{}first\PYGZus{}div\PYGZus{}freq(uint32\PYGZus{}t target\PYGZus{}freq, hal\PYGZus{}clk\PYGZus{}t next\PYGZus{}clk)
 \PYGZob{}
        int ret = 0;
@@ \PYGZhy{}339,3 +378,20 @@ int cpufreq\PYGZus{}vf\PYGZus{}init(void)

        return 0;
 \PYGZcb{}
\end{sphinxVerbatim}

\sphinxAtStartPar
测试：

\sphinxAtStartPar
设置LDO电压为1.08V：

\begin{sphinxVerbatim}[commandchars=\\\{\}]
c906\PYGZgt{}set\PYGZus{}cpu\PYGZus{}voltage 1080
set\PYGZus{}cpu\PYGZus{}voltage ok, voltage: 1080 mV

c906\PYGZgt{}get\PYGZus{}cpu\PYGZus{}voltage
get\PYGZus{}cpu\PYGZus{}voltage ok, voltage: 1100 mV 
\end{sphinxVerbatim}

\sphinxAtStartPar
设置LDO电压1.08V成功，但是最后读取出来的LDO电压为1.1V。

\sphinxAtStartPar
原因：

\sphinxAtStartPar
根据宏定义\#define VOLTAGE\_STEP 25以及set\_cpu\_voltage函数的具体实现：

\begin{sphinxVerbatim}[commandchars=\\\{\}]
ldo\PYGZus{}vol\PYGZus{}field = ldo\PYGZus{}cal + ((int)vol \PYGZhy{} LDO\PYGZus{}CAL\PYGZus{}BASE) / VOLTAGE\PYGZus{}STEP;
\end{sphinxVerbatim}

\sphinxAtStartPar
可以知道步进电压为25mV，也就是设置的电压必须为25mV的整数倍，因此最接近1.08V的电压为1.075V，可以尝试设置LDO电压值为1075。

\begin{sphinxVerbatim}[commandchars=\\\{\}]
c906\PYGZgt{}set\PYGZus{}cpu\PYGZus{}voltage 1075
set\PYGZus{}cpu\PYGZus{}voltage ok, voltage: 1075 mV

c906\PYGZgt{}get\PYGZus{}cpu\PYGZus{}voltage
get\PYGZus{}cpu\PYGZus{}voltage ok, voltage: 1075 mV
\end{sphinxVerbatim}

\sphinxAtStartPar
可以看到成功设置为LDO电压为1.075V。


\chapter{验证}
\label{\detokenize{r128_u9a8c_u8bc1_u56fa_u4ef6/r128_u9a8c_u8bc1_u56fa_u4ef6_u5f00_u53d1_u6307_u5357:id1}}

\section{1、 M33、DSP 均以默认频率运行，以 DPLL1可配置频率为主 ，默认频率为M33\sphinxhyphen{}192M，RV\sphinxhyphen{}533M，DSP\sphinxhyphen{}400M；}
\label{\detokenize{r128_u9a8c_u8bc1_u56fa_u4ef6/r128_u9a8c_u8bc1_u56fa_u4ef6_u5f00_u53d1_u6307_u5357:id2}}
\sphinxAtStartPar
\sphinxstylestrong{M33固件下，重启查看log：}

\begin{sphinxVerbatim}[commandchars=\\\{\}]
M33 CPU Clock Freq: 192 MHz
C906 CPU Freq: 480 MHz
DSP core clock frequency: 384000000Hz
\end{sphinxVerbatim}

\sphinxAtStartPar
\sphinxstylestrong{RV固件下，重启查看log：}

\begin{sphinxVerbatim}[commandchars=\\\{\}]
M33 CPU Clock Freq: 192 MHz
C906 CPU Freq: 533 MHz
DSP core clock frequency: 384000000Hz
\end{sphinxVerbatim}


\section{2、将HS\sphinxhyphen{}PSRAM 从DPLL3切换到DPLL1中 ，频率根据当前条件 尽量配高频 ， 默认运行；}
\label{\detokenize{r128_u9a8c_u8bc1_u56fa_u4ef6/r128_u9a8c_u8bc1_u56fa_u4ef6_u5f00_u53d1_u6307_u5357:hs-psram-dpll3dpll1}}
\sphinxAtStartPar
\sphinxstylestrong{在RV或者M33固件下，执行以下命令：}

\begin{sphinxVerbatim}[commandchars=\\\{\}]
c906\PYGZgt{}p 0x4004c4e0 //读取system clock control寄存器
reg\PYGZus{}addr[0x4004c4e0]=0x00022191  //bit21为0，说明HS\PYGZhy{}PSRAM时钟源为DPLL1
\end{sphinxVerbatim}

\sphinxAtStartPar
\sphinxincludegraphics{{CKPLL_HSPSRAM_SEL}.png}


\section{3.1、将 RV挂载到DPLL3 上，DPLL3的时钟源尽量只给RV使用；}
\label{\detokenize{r128_u9a8c_u8bc1_u56fa_u4ef6/r128_u9a8c_u8bc1_u56fa_u4ef6_u5f00_u53d1_u6307_u5357:id3}}
\sphinxAtStartPar
\sphinxstylestrong{在RV固件下，执行以下命令：}

\begin{sphinxVerbatim}[commandchars=\\\{\}]
c906\PYGZgt{}p 0x4004c4e0 //读取system clock control寄存器
reg\PYGZus{}addr[0x4004c4e0]=0x00022191  //bit17为1，说明RV时钟源为DPLL3
\end{sphinxVerbatim}

\sphinxAtStartPar
\sphinxincludegraphics{{CKPLL_C906_SEL}.png}

\sphinxAtStartPar
bit18\textasciitilde{}bit21为0，表示HSPSRAM、LSPSRAM、M33、DSP的时钟源都是DPLL1，DPLL3的时钟源只给RV使用。


\section{3.2、将 M33挂载到DPLL3 上，DPLL3的时钟源尽量只给M33使用；}
\label{\detokenize{r128_u9a8c_u8bc1_u56fa_u4ef6/r128_u9a8c_u8bc1_u56fa_u4ef6_u5f00_u53d1_u6307_u5357:m33dpll3-dpll3m33}}
\sphinxAtStartPar
\sphinxstylestrong{在M33固件下，执行以下命令：}

\begin{sphinxVerbatim}[commandchars=\\\{\}]
c906\PYGZgt{}p 0x4004c4e0//读取system clock control寄存器
reg\PYGZus{}addr[0x4004c4e0]=0x00082191//bit19为1，说明M33时钟源为DPLL3
\end{sphinxVerbatim}

\sphinxAtStartPar
\sphinxincludegraphics{{CKPLL_M33_SEL}.png}

\sphinxAtStartPar
bit18、bit20和bit21为0，表示HSPSRAM、LSPSRAM、RV、DSP的时钟源都是DPLL1，DPLL3的时钟源只给M33使用。


\section{4、支持 RV跑memtester ，跑在HS\sphinxhyphen{}PSRAM上；}
\label{\detokenize{r128_u9a8c_u8bc1_u56fa_u4ef6/r128_u9a8c_u8bc1_u56fa_u4ef6_u5f00_u53d1_u6307_u5357:id4}}
\sphinxAtStartPar
\sphinxstylestrong{在RV或者M33固件下，执行以下命令：}

\begin{sphinxVerbatim}[commandchars=\\\{\}]
c906\PYGZgt{}memtester \PYGZhy{}s 8k \PYGZhy{}t 1
memtester version 4.3.0 (64\PYGZhy{}bit)
Copyright (C) 2001\PYGZhy{}2012 Charles Cazabon.
Licensed under the GNU General Public License version 2 (only).

sysconf(\PYGZus{}SC\PYGZus{}PAGE\PYGZus{}SIZE) not supported; using pagesize of 8192
memsuffix k
want 0MB (8192 bytes)
got  0MB (8192 bytes) in 0xc643d20Loop 1/1:
\PYGZlt{}============Prepare============\PYGZgt{}
  Stuck Address       : ok
  Random Value        :  ok
\PYGZlt{}============Start test============\PYGZgt{}
  Compare XOR         : ok
  Compare SUB         : ok
  Compare MUL         : ok
  Compare DIV         : ok
  Compare OR          : ok
  Compare AND         : ok
  Sequential Increment: ok
  Solid Bits          : ok
  Block Sequential    : ok
  Checkerboard        : ok
  Bit Spread          : ok
  Bit Flip            : ok
  Walking Ones        : ok
  Walking Zeroes      : ok

Done.

c906\PYGZgt{}rpccli arm memtester \PYGZhy{}s 8k \PYGZhy{}t 1 
memtester version 4.3.0 (32\PYGZhy{}bit)
Copyright (C) 2001\PYGZhy{}2012 Charles Cazabon.
Licensed under the GNU General Public License version 2 (only).

sysconf(\PYGZus{}SC\PYGZus{}PAGE\PYGZus{}SIZE) not supported; using pagesize of 8192
memsuffix k
want 0MB (8192 bytes)
got  0MB (8192 bytes) in 0xa5a5a5a50c13c220Loop 1/1:
\PYGZlt{}============Prepare============\PYGZgt{}
  Stuck Address       : ok
  Random Value        :  ok
\PYGZlt{}============Start test============\PYGZgt{}
  Compare XOR         : ok
  Compare SUB         : ok
  Compare MUL         : ok
  Compare DIV         : ok
  Compare OR          : ok
  Compare AND         : ok
  Sequential Increment: ok
  Solid Bits          : ok
  Block Sequential    : ok
  Checkerboard        : ok
  Bit Spread          : ok
  Bit Flip            : ok
  Walking Ones        : ok
  Walking Zeroes      : ok

Done.
\end{sphinxVerbatim}

\sphinxAtStartPar
查看RV核和M33核的启动地址：

\sphinxAtStartPar
\sphinxincludegraphics{{RISCV_CPU_Start_Address}.png}

\sphinxAtStartPar
地址为0xC开头，表明RV和M33都运行在RPSRAM中。


\section{5、支持查看RV占用率，确保RV运行memtester正常；}
\label{\detokenize{r128_u9a8c_u8bc1_u56fa_u4ef6/r128_u9a8c_u8bc1_u56fa_u4ef6_u5f00_u53d1_u6307_u5357:id5}}
\sphinxAtStartPar
\sphinxstylestrong{RV或M33固件下，直接运行run\sphinxhyphen{}time\sphinxhyphen{}stats}

\begin{sphinxVerbatim}[commandchars=\\\{\}]
c906\PYGZgt{}run\PYGZhy{}time\PYGZhy{}stats
Task            Abs Time      \PYGZpc{} Time
****************************************
CLI             616             \PYGZlt{}1\PYGZpc{}
IDLE            516619          99\PYGZpc{}
adbd\PYGZhy{}input      0               \PYGZlt{}1\PYGZpc{}
tcpip           0               \PYGZlt{}1\PYGZpc{}
nand\PYGZus{}rcd        0               \PYGZlt{}1\PYGZpc{}
nftld           0               \PYGZlt{}1\PYGZpc{}
amp\PYGZhy{}admin       0               \PYGZlt{}1\PYGZpc{}
Tmr Svc         0               \PYGZlt{}1\PYGZpc{}
amp\PYGZhy{}ser1        0               \PYGZlt{}1\PYGZpc{}
amp\PYGZhy{}ser2        0               \PYGZlt{}1\PYGZpc{}
amp\PYGZhy{}ser3        0               \PYGZlt{}1\PYGZpc{}
amp\PYGZhy{}ser4        0               \PYGZlt{}1\PYGZpc{}
pm\PYGZus{}suspend      0               \PYGZlt{}1\PYGZpc{}
hub\PYGZhy{}main\PYGZhy{}thread 0               \PYGZlt{}1\PYGZpc{}
amp\PYGZhy{}send\PYGZhy{}task   0               \PYGZlt{}1\PYGZpc{}
amp\PYGZhy{}recv\PYGZhy{}task   0               \PYGZlt{}1\PYGZpc{}
udc\PYGZus{}vbus\PYGZus{}det    0               \PYGZlt{}1\PYGZpc{}
adbd\PYGZhy{}output     1               \PYGZlt{}1\PYGZpc{}
AudioMT2        0               \PYGZlt{}1\PYGZpc{}
cpu\PYGZus{}zone        0               \PYGZlt{}1\PYGZpc{}
disp2           258             \PYGZlt{}1\PYGZpc{}
amp\PYGZhy{}ser0        0               \PYGZlt{}1\PYGZpc{}
\end{sphinxVerbatim}

\sphinxAtStartPar
6、 将RV、M33内部LDO供电电压修改至1.08V，DSP内部LDO供电保持1.175V不变；

\sphinxAtStartPar
\sphinxstylestrong{RV固件下：}

\begin{sphinxVerbatim}[commandchars=\\\{\}]
c906\PYGZgt{}set\PYGZus{}cpu\PYGZus{}voltage 1075 //设置RV内部LDO供电电压
set\PYGZus{}cpu\PYGZus{}voltage ok, voltage: 1075 mV

c906\PYGZgt{}get\PYGZus{}cpu\PYGZus{}voltage  //读取RV内部LDO供电电压
get\PYGZus{}cpu\PYGZus{}voltage ok, voltage: 1075 mV 
\end{sphinxVerbatim}

\sphinxAtStartPar
\sphinxstylestrong{M33固件下：}

\begin{sphinxVerbatim}[commandchars=\\\{\}]
c906\PYGZgt{}rpccli arm set\PYGZus{}cpu\PYGZus{}voltage 1075 //设置M33内部LDO供电电压
set\PYGZus{}cpu\PYGZus{}voltage ok, voltage: 1075 mV

c906\PYGZgt{}rpccli arm get\PYGZus{}cpu\PYGZus{}voltage //读取M33内部LDO供电电压
get\PYGZus{}cpu\PYGZus{}voltage ok, voltage: 1075 mV
\end{sphinxVerbatim}


\chapter{M33验证固件}
\label{\detokenize{r128_u9a8c_u8bc1_u56fa_u4ef6/r128_u9a8c_u8bc1_u56fa_u4ef6_u5f00_u53d1_u6307_u5357:m33}}

\section{将M33挂载到DPLL3上，DPLL3的时钟源尽量只给M33使用}
\label{\detokenize{r128_u9a8c_u8bc1_u56fa_u4ef6/r128_u9a8c_u8bc1_u56fa_u4ef6_u5f00_u53d1_u6307_u5357:id6}}
\sphinxAtStartPar
相关源代码文件：\sphinxcode{\sphinxupquote{r128\sphinxhyphen{}dev/lichee/rtos/arch/arm/armv8m/sun20iw2p1/sun20i.c}}

\sphinxAtStartPar
相关函数：\sphinxcode{\sphinxupquote{static int ar200a\_clk\_set(u32 freq)}}

\sphinxAtStartPar
相关代码补丁：

\begin{sphinxVerbatim}[commandchars=\\\{\}]
diff \PYGZhy{}\PYGZhy{}git a/arch/arm/armv8m/sun20iw2p1/sun20i.c b/arch/arm/armv8m/sun20iw2p1/sun20i.c
index 1c7cb1abd..509f3b516 100755
\PYGZhy{}\PYGZhy{}\PYGZhy{} a/arch/arm/armv8m/sun20iw2p1/sun20i.c
+++ b/arch/arm/armv8m/sun20iw2p1/sun20i.c
@@ \PYGZhy{}173,12 +173,16 @@ err:
  * we fixed ck1\PYGZus{}m33 to 384M,
  * so freq only support 384M,192M,128M,96M,64M,48M,32M,24M
  */
+ \PYGZsh{}define M33\PYGZus{}DPLL1\PYGZus{}CLK 1
+ \PYGZsh{}define M33\PYGZus{}DPLL3\PYGZus{}CLK 3
+ \PYGZsh{}define M33\PYGZus{}PLL\PYGZus{}CLK\PYGZus{}SOURCE M33\PYGZus{}DPLL3\PYGZus{}CLK
 static int ar200a\PYGZus{}clk\PYGZus{}set(u32 freq)
 \PYGZob{}
        u32 rate;
\PYGZhy{}       hal\PYGZus{}clk\PYGZus{}t clk\PYGZus{}ck1\PYGZus{}m33, clk\PYGZus{}ck\PYGZus{}m33, clk\PYGZus{}sys, clk\PYGZus{}ar200a\PYGZus{}f;
+       hal\PYGZus{}clk\PYGZus{}t clk\PYGZus{}ck1\PYGZus{}m33, clk\PYGZus{}ck3\PYGZus{}m33, clk\PYGZus{}ck\PYGZus{}m33, clk\PYGZus{}sys, clk\PYGZus{}ar200a\PYGZus{}f;
        hal\PYGZus{}clk\PYGZus{}status\PYGZus{}t ret = 0;

+\PYGZsh{}if (M33\PYGZus{}PLL\PYGZus{}CLK\PYGZus{}SOURCE == M33\PYGZus{}DPLL1\PYGZus{}CLK)
        /*set ck1\PYGZus{}m33\PYGZus{}clk*/
        //sr32(CCMU\PYGZus{}AON\PYGZus{}BASE+0xa4,  3, 1, 1);
        //0x4004c4a4: 0x80800008
@@ \PYGZhy{}207,12 +211,25 @@ static int ar200a\PYGZus{}clk\PYGZus{}set(u32 freq)
                ret = \PYGZhy{}1;
                goto err3;
        \PYGZcb{}
\PYGZhy{}
+       printf(\PYGZdq{}M33: chose clock source to DPLL1\PYGZbs{}n\PYGZdq{});
+\PYGZsh{}else
+       printf(\PYGZdq{}DPLL3\PYGZbs{}n\PYGZdq{});
+       /*enable ck3\PYGZus{}m33\PYGZus{}clk*/
+       sr32(CCMU\PYGZus{}AON\PYGZus{}BASE+0xa8,  3, 1, 1);
+       //0x4004c4a4: 0x80800008
+       /*fixed 384M*/
+       sr32(CCMU\PYGZus{}AON\PYGZus{}BASE+0xa8,  0, 3, 0x3);
+       //0x4004c4a4: 0x8080000b
+       /*set ck\PYGZus{}m33\PYGZus{}clk*/
+       sr32(CCMU\PYGZus{}AON\PYGZus{}BASE+0xe0, 19, 1, 1);
+       //0x4004c4e0: 0x00000101
+       printf(\PYGZdq{}M33: chose clock source to DPLL3\PYGZbs{}n\PYGZdq{});
+\PYGZsh{}endif
        /*set sys\PYGZus{}clk*/
        //sr32(CCMU\PYGZus{}AON\PYGZus{}BASE+0xe0,  8, 4, 0x7);
        //0x4004c4e0: 0x00000701
        clk\PYGZus{}sys = hal\PYGZus{}clock\PYGZus{}get(HAL\PYGZus{}SUNXI\PYGZus{}AON\PYGZus{}CCU, CLK\PYGZus{}SYS);
\PYGZhy{}       ret = hal\PYGZus{}clk\PYGZus{}set\PYGZus{}rate(clk\PYGZus{}sys, freq);
+       ret = hal\PYGZus{}clk\PYGZus{}set\PYGZus{}rate(clk\PYGZus{}sys, freq); //freq=192MHz
        if (HAL\PYGZus{}CLK\PYGZus{}STATUS\PYGZus{}OK != ret) \PYGZob{}
                ret = \PYGZhy{}1;
                goto err4;
...
\end{sphinxVerbatim}


\chapter{查看cpu频率}
\label{\detokenize{r128_u9a8c_u8bc1_u56fa_u4ef6/r128_u9a8c_u8bc1_u56fa_u4ef6_u5f00_u53d1_u6307_u5357:cpu}}

\section{M33和C906}
\label{\detokenize{r128_u9a8c_u8bc1_u56fa_u4ef6/r128_u9a8c_u8bc1_u56fa_u4ef6_u5f00_u53d1_u6307_u5357:m33c906}}
\sphinxAtStartPar
\sphinxstylestrong{（1）确定HOSC时钟源：}

\begin{sphinxVerbatim}[commandchars=\\\{\}]
p 0x4004c484 //查看bit0:2的值， 根据下面信息直接确定f\PYGZus{}hosc的值
\end{sphinxVerbatim}

\sphinxAtStartPar
\sphinxincludegraphics{{hosc_reg}.png}

\sphinxAtStartPar
\sphinxstylestrong{（2）确定DPLL3的时钟频率：}

\begin{sphinxVerbatim}[commandchars=\\\{\}]
p 0x4004c494 //f\PYGZus{}dpll3=(f\PYGZus{}hosc*N)/M，N由bit4:11决定，M由0:3决定
\end{sphinxVerbatim}

\sphinxAtStartPar
\sphinxincludegraphics{{DPLL3_REG}.png}

\sphinxAtStartPar
\sphinxstylestrong{（3）确定M33的频率}

\sphinxAtStartPar
\sphinxincludegraphics{{M33_freq}.png}

\begin{sphinxVerbatim}[commandchars=\\\{\}]
c906\PYGZgt{}p 0x4004c4a8 //查看bit0:2的值，直接确定分频系数M，那么f\PYGZus{}m33=f\PYGZus{}dpll3/M
\end{sphinxVerbatim}

\sphinxAtStartPar
\sphinxstylestrong{（4）确定C906的频率}

\sphinxAtStartPar
\sphinxincludegraphics{{C906_freq}.png}

\begin{sphinxVerbatim}[commandchars=\\\{\}]
c906\PYGZgt{}p 0x4004c4a8 //查看bit4:6的值，直接确定分频系数M，那么f\PYGZus{}c906=f\PYGZus{}dpll3/M
\end{sphinxVerbatim}


\section{HPSRAM}
\label{\detokenize{r128_u9a8c_u8bc1_u56fa_u4ef6/r128_u9a8c_u8bc1_u56fa_u4ef6_u5f00_u53d1_u6307_u5357:hpsram}}
\sphinxAtStartPar
\sphinxincludegraphics{{CK1_HPSRAM_REG}.png}

\begin{sphinxVerbatim}[commandchars=\\\{\}]
p 0x4004c4a8 //查看分频系数M，直接计算f\PYGZus{}hpsram=f\PYGZus{}dpll1/M
\end{sphinxVerbatim}

\sphinxstepscope


\chapter{MyST Markdown 语法入门}
\label{\detokenize{other/myst_markdown:myst-markdown}}\label{\detokenize{other/myst_markdown::doc}}
\sphinxAtStartPar
MyST Markdown, MyST（全称 Markedly Structured Text）专为 Sphinx 项目使用而设计的 Markdown 风格。 它是 CommonMark Markdown 和一些额外的语法片段的组合，以支持 Sphinx 的功能，因此您可以用纯 Markdown 编写 Sphinx 文档。

\sphinxAtStartPar
\sphinxhref{https://myst-parser.readthedocs.io}{MyST\sphinxhyphen{}Parser官网语法教程}


\section{提示框}
\label{\detokenize{other/myst_markdown:id1}}
\sphinxAtStartPar
警告提示框（Admonition）突出显示与页面叙述稍有不同的特定文本块，例如注释或警告。

\begin{sphinxVerbatim}[commandchars=\\\{\}]
:::\PYGZob{}tip\PYGZcb{}
给读者一些有用的提示吧！
:::

:::\PYGZob{}note\PYGZcb{}
这是一个笔记
:::


:::\PYGZob{}caution\PYGZcb{}
这是一个警告
:::

:::\PYGZob{}error\PYGZcb{}
这是一个错误
:::
\end{sphinxVerbatim}

\sphinxAtStartPar
显示效果

\sphinxAtStartPar
:::\{tip\}
给读者一些有用的提示吧！
:::

\sphinxAtStartPar
:::\{note\}
这是一个笔记
:::

\sphinxAtStartPar
:::\{caution\}
这是一个警告
:::

\sphinxAtStartPar
:::\{error\}
这是一个错误
:::


\bigskip\hrule\bigskip



\section{标题}
\label{\detokenize{other/myst_markdown:id2}}
\sphinxAtStartPar
为了兼容考虑，用一个空格在 \sphinxcode{\sphinxupquote{\#}} 和标题之间进行分隔。

\begin{sphinxVerbatim}[commandchars=\\\{\}]
\PYG{c+c1}{\PYGZsh{} 一级标题}

\PYG{c+c1}{\PYGZsh{}\PYGZsh{} 二级标题}

\PYG{c+c1}{\PYGZsh{}\PYGZsh{}\PYGZsh{} 三级标题}

\PYG{c+c1}{\PYGZsh{}\PYGZsh{}\PYGZsh{}\PYGZsh{} 四级标题}

\PYG{c+c1}{\PYGZsh{}\PYGZsh{}\PYGZsh{}\PYGZsh{}\PYGZsh{} 五级标题}
\end{sphinxVerbatim}


\bigskip\hrule\bigskip



\section{有序列表}
\label{\detokenize{other/myst_markdown:id3}}
\sphinxAtStartPar
要创建有序列表，请在每个列表项前添加数字并紧跟一个英文句点。\\
数字不必按数学顺序排列，但是列表应当以数字 1 起始。

\begin{sphinxVerbatim}[commandchars=\\\{\}]
\PYG{k}{1.} 条目1
\PYG{k}{2.} 条目2
\PYG{k}{3.} 条目3
\PYG{+w}{    }\PYG{k}{\PYGZhy{}}\PYG{+w}{ }缩进条目
\PYG{+w}{    }\PYG{k}{\PYGZhy{}}\PYG{+w}{ }缩进条目
\end{sphinxVerbatim}

\sphinxAtStartPar
显示效果
\begin{enumerate}
\sphinxsetlistlabels{\arabic}{enumi}{enumii}{}{.}%
\item {} 
\sphinxAtStartPar
条目1

\item {} 
\sphinxAtStartPar
条目2

\item {} 
\sphinxAtStartPar
条目3
\begin{itemize}
\item {} 
\sphinxAtStartPar
缩进条目

\item {} 
\sphinxAtStartPar
缩进条目

\end{itemize}

\end{enumerate}


\bigskip\hrule\bigskip



\section{无序列表}
\label{\detokenize{other/myst_markdown:id4}}
\sphinxAtStartPar
要创建无序列表，请在每个列表项前面添加破折号 (\sphinxhyphen{})、星号 (*) 或加号 (+) 。缩进一个或多个列表项可创建嵌套列表。

\begin{sphinxVerbatim}[commandchars=\\\{\}]
\PYG{k}{\PYGZhy{}}\PYG{+w}{ }条目1
\PYG{k}{\PYGZhy{}}\PYG{+w}{ }条目2
\PYG{k}{\PYGZhy{}}\PYG{+w}{ }条目3
\PYG{k}{\PYGZhy{}}\PYG{+w}{ }条目4
\PYG{+w}{    }\PYG{k}{\PYGZhy{}}\PYG{+w}{ }缩进条目
\PYG{+w}{    }\PYG{k}{\PYGZhy{}}\PYG{+w}{ }缩进条目

\PYG{k}{*}\PYG{+w}{ }条目1
\PYG{k}{*}\PYG{+w}{ }条目2
\PYG{k}{*}\PYG{+w}{ }条目3
\PYG{k}{*}\PYG{+w}{ }条目4
\PYG{+w}{    }\PYG{k}{*}\PYG{+w}{ }缩进条目
\PYG{+w}{    }\PYG{k}{*}\PYG{+w}{ }缩进条目

+ 条目1
+ 条目2
+ 条目3
+ 条目4
    + 缩进条目
    + 缩进条目
\end{sphinxVerbatim}

\sphinxAtStartPar
显示效果
\begin{itemize}
\item {} 
\sphinxAtStartPar
条目1

\item {} 
\sphinxAtStartPar
条目2

\item {} 
\sphinxAtStartPar
条目3

\item {} 
\sphinxAtStartPar
条目4
\begin{itemize}
\item {} 
\sphinxAtStartPar
缩进条目

\item {} 
\sphinxAtStartPar
缩进条目

\end{itemize}

\end{itemize}
\begin{itemize}
\item {} 
\sphinxAtStartPar
条目1

\item {} 
\sphinxAtStartPar
条目2

\item {} 
\sphinxAtStartPar
条目3

\item {} 
\sphinxAtStartPar
条目4
\begin{itemize}
\item {} 
\sphinxAtStartPar
缩进条目

\item {} 
\sphinxAtStartPar
缩进条目

\end{itemize}

\end{itemize}
\begin{itemize}
\item {} 
\sphinxAtStartPar
条目1

\item {} 
\sphinxAtStartPar
条目2

\item {} 
\sphinxAtStartPar
条目3

\item {} 
\sphinxAtStartPar
条目4
\begin{itemize}
\item {} 
\sphinxAtStartPar
缩进条目

\item {} 
\sphinxAtStartPar
缩进条目

\end{itemize}

\end{itemize}


\bigskip\hrule\bigskip



\section{文本换行}
\label{\detokenize{other/myst_markdown:id5}}
\sphinxAtStartPar
\sphinxstylestrong{Markdown中在段落前不要用空格（spaces）或制表符（ tabs）缩进段落。}


\subsection{无效换行}
\label{\detokenize{other/myst_markdown:id6}}
\sphinxAtStartPar
在一行结尾没有加空格时，后一行内容会自动追加显示在前一行

\begin{sphinxVerbatim}[commandchars=\\\{\}]
This is the first line.  
And this is the second line.
\end{sphinxVerbatim}

\sphinxAtStartPar
显示效果

\sphinxAtStartPar
This is the first line.
And this is the second line.


\subsection{段内换行}
\label{\detokenize{other/myst_markdown:id7}}
\sphinxAtStartPar
要在没有段落的情况下换行，使用\sphinxstylestrong{两个或多个空格}进行换行，称为 结尾空格（trailing whitespace)

\begin{sphinxVerbatim}[commandchars=\\\{\}]
This is the first line.  
And this is the second line.
\end{sphinxVerbatim}

\sphinxAtStartPar
显示效果

\sphinxAtStartPar
This is the first line.

\sphinxAtStartPar
And this is the second line.


\subsection{新段落}
\label{\detokenize{other/myst_markdown:id8}}
\sphinxAtStartPar
两行之间\sphinxstylestrong{间隔一行}，即可新增段落

\begin{sphinxVerbatim}[commandchars=\\\{\}]
This is the first line.

And this is the second line.
\end{sphinxVerbatim}

\sphinxAtStartPar
显示效果

\sphinxAtStartPar
This is the first line.

\sphinxAtStartPar
And this is the second line.


\bigskip\hrule\bigskip



\section{字体样式}
\label{\detokenize{other/myst_markdown:id9}}
\sphinxAtStartPar
要加粗文本，在短语前后各添加\sphinxstylestrong{两个}星号\sphinxcode{\sphinxupquote{*}}或下划线\sphinxcode{\sphinxupquote{\_}}\\
要斜体文本，在短语前后各添加\sphinxstylestrong{一个}星号\sphinxcode{\sphinxupquote{*}}或下划线\sphinxcode{\sphinxupquote{\_}}\\
要同时用粗体和斜体突出显示文本，在短语前后各添加\sphinxstylestrong{三个}星号\sphinxcode{\sphinxupquote{*}}或下划线\sphinxcode{\sphinxupquote{\_}}\\
删除线在PDF中不可用

\begin{sphinxVerbatim}[commandchars=\\\{\}]
\PYG{g+ge}{*斜体*} 

\PYG{g+ge}{\PYGZus{}斜体\PYGZus{}}

\PYG{g+gs}{**粗体**}

\PYG{g+gs}{\PYGZus{}\PYGZus{}粗体\PYGZus{}\PYGZus{}}

***粗体+斜体***

\PYGZus{}\PYGZus{}\PYGZus{}粗体+斜体\PYGZus{}\PYGZus{}\PYGZus{}


H\PYGZob{}sub\PYGZcb{}`2`O, and 4\PYGZob{}sup\PYGZcb{}`th` of July
\end{sphinxVerbatim}

\sphinxAtStartPar
显示效果

\sphinxAtStartPar
\sphinxstyleemphasis{斜体}

\sphinxAtStartPar
\sphinxstyleemphasis{斜体}

\sphinxAtStartPar
\sphinxstylestrong{粗体}

\sphinxAtStartPar
\sphinxstylestrong{粗体}

\sphinxAtStartPar
\sphinxstyleemphasis{\sphinxstylestrong{粗体+斜体}}

\sphinxAtStartPar
\sphinxstyleemphasis{\sphinxstylestrong{粗体+斜体}}

\sphinxAtStartPar
H$_{\text{2}}$O, and 4$^{\text{th}}$ of July


\bigskip\hrule\bigskip



\section{单行块引用}
\label{\detokenize{other/myst_markdown:id10}}
\sphinxAtStartPar
要创建块引用，请在段落前添加一个\sphinxcode{\sphinxupquote{>}}符号。

\begin{sphinxVerbatim}[commandchars=\\\{\}]
\PYG{k}{\PYGZgt{} }\PYG{g+ge}{这是一个引用块}
\end{sphinxVerbatim}

\sphinxAtStartPar
显示效果
\begin{quote}

\sphinxAtStartPar
这是一个引用块
\end{quote}


\bigskip\hrule\bigskip



\section{多行块引用}
\label{\detokenize{other/myst_markdown:id11}}
\sphinxAtStartPar
块引用可以包含多个段落。为段落之间的空白行添加一个 \sphinxcode{\sphinxupquote{>}}符号。

\begin{sphinxVerbatim}[commandchars=\\\{\}]
\PYG{k}{\PYGZgt{} }\PYG{g+ge}{呼啸的秋风让人无限忧愁，进也忧愁，退也忧愁。}
\PYGZgt{} 
\PYG{k}{\PYGZgt{} }\PYG{g+ge}{异域戍边的人，哪个不陷入悲愁中？真是愁白了头啊。}
\PYGZgt{} 
\PYG{k}{\PYGZgt{} }\PYG{g+ge}{胡人之处多狂风，树木萧瑟干枯。}
\end{sphinxVerbatim}

\sphinxAtStartPar
显示效果
\begin{quote}

\sphinxAtStartPar
呼啸的秋风让人无限忧愁，进也忧愁，退也忧愁。

\sphinxAtStartPar
异域戍边的人，哪个不陷入悲愁中？真是愁白了头啊。

\sphinxAtStartPar
胡人之处多狂风，树木萧瑟干枯。
\end{quote}


\bigskip\hrule\bigskip



\section{代码块}
\label{\detokenize{other/myst_markdown:id12}}
\sphinxAtStartPar
在代码前和后的行上使用三个反引号\sphinxcode{\sphinxupquote{```}}或三个波浪号\sphinxcode{\sphinxupquote{\textasciitilde{}\textasciitilde{}\textasciitilde{}}}
在代码块之前的反引号旁边指定一种语言，即可实现\sphinxstylestrong{语法高亮}

\begin{sphinxVerbatim}[commandchars=\\\{\}]
\PYG{l+s+sb}{    ```}\PYG{l+s+sb}{python}
    \PYG{k+kn}{import}\PYG{+w}{ }\PYG{n+nn}{os}
    \PYG{n+nb}{print}\PYG{p}{(}\PYG{l+s+s2}{\PYGZdq{}}\PYG{l+s+s2}{Hello World!}\PYG{l+s+s2}{\PYGZdq{}}\PYG{p}{)}
\PYG{l+s+sb}{    ```}

    \PYGZti{}\PYGZti{}\PYGZti{} json
    \PYGZob{}
        \PYGZdq{}firstName\PYGZdq{}: \PYGZdq{}John\PYGZdq{},
        \PYGZdq{}lastName\PYGZdq{}: \PYGZdq{}Smith\PYGZdq{},
        \PYGZdq{}age\PYGZdq{}: 25
    \PYGZcb{}
    \PYGZti{}\PYGZti{}\PYGZti{}
\end{sphinxVerbatim}

\sphinxAtStartPar
显示效果

\begin{sphinxVerbatim}[commandchars=\\\{\}]
\PYG{k+kn}{import}\PYG{+w}{ }\PYG{n+nn}{os}
\PYG{n+nb}{print}\PYG{p}{(}\PYG{l+s+s2}{\PYGZdq{}}\PYG{l+s+s2}{Hello World!}\PYG{l+s+s2}{\PYGZdq{}}\PYG{p}{)}
\end{sphinxVerbatim}

\begin{sphinxVerbatim}[commandchars=\\\{\}]
\PYG{p}{\PYGZob{}}
\PYG{+w}{    }\PYG{n+nt}{\PYGZdq{}firstName\PYGZdq{}}\PYG{p}{:}\PYG{+w}{ }\PYG{l+s+s2}{\PYGZdq{}John\PYGZdq{}}\PYG{p}{,}
\PYG{+w}{    }\PYG{n+nt}{\PYGZdq{}lastName\PYGZdq{}}\PYG{p}{:}\PYG{+w}{ }\PYG{l+s+s2}{\PYGZdq{}Smith\PYGZdq{}}\PYG{p}{,}
\PYG{+w}{    }\PYG{n+nt}{\PYGZdq{}age\PYGZdq{}}\PYG{p}{:}\PYG{+w}{ }\PYG{l+m+mi}{25}
\PYG{p}{\PYGZcb{}}
\end{sphinxVerbatim}


\bigskip\hrule\bigskip



\section{文件中的代码}
\label{\detokenize{other/myst_markdown:id13}}
\sphinxAtStartPar
可以使用 literalinclude 指令从文件中包含更长的代码：\\
\sphinxstylestrong{注意:} ```和\{literalinclude\}之间不能有空格

\begin{sphinxVerbatim}[commandchars=\\\{\}]
    ```\PYGZob{}literalinclude\PYGZcb{} quickstart.py
    ```
\end{sphinxVerbatim}


\bigskip\hrule\bigskip



\section{数学公式}
\label{\detokenize{other/myst_markdown:id14}}
\sphinxAtStartPar
math角色和指令分别用于定义内联数学和块数学。\\
\sphinxstylestrong{注意:} ```和\{math\}之间不能有空格

\begin{sphinxVerbatim}[commandchars=\\\{\}]
    ```\PYGZob{}math\PYGZcb{}
    :label: mymath
    (a + b)\PYGZca{}2 = a\PYGZca{}2 + 2ab + b\PYGZca{}2

    (a + b)\PYGZca{}2  \PYGZam{}=  (a + b)(a + b) \PYGZbs{}\PYGZbs{}
            \PYGZam{}=  a\PYGZca{}2 + 2ab + b\PYGZca{}2
    ```
\end{sphinxVerbatim}

\sphinxAtStartPar
显示效果
\begin{align}\label{equation:other/myst_markdown:mymath}\!\begin{aligned}
(a + b)^2 = a^2 + 2ab + b^2\\
(a + b)^2  &=  (a + b)(a + b) \\
           &=  a^2 + 2ab + b^2\\
\end{aligned}\end{align}

\bigskip\hrule\bigskip



\section{脚注}
\label{\detokenize{other/myst_markdown:id15}}
\sphinxAtStartPar
脚注使用 pandoc 规范。它们的标签以 \textasciicircum{} 开头，然后可以是任何字母数字字符串（无空格），不区分大小写。

\begin{sphinxVerbatim}[commandchars=\\\{\}]
\PYG{k}{\PYGZhy{}}\PYG{+w}{ }This is a manually\PYGZhy{}numbered footnote reference.[\PYGZca{}3]
\PYG{k}{\PYGZhy{}}\PYG{+w}{ }This is an auto\PYGZhy{}numbered footnote reference.[\PYGZca{}myref]

[\PYG{n+nl}{\PYGZca{}myref}]: \PYG{n+na}{This is an auto\PYGZhy{}numbered footnote definition.}
[\PYG{n+nl}{\PYGZca{}3}]: \PYG{n+na}{This is a manually\PYGZhy{}numbered footnote definition.}
\end{sphinxVerbatim}

\sphinxAtStartPar
显示效果
\begin{itemize}
\item {} 
\sphinxAtStartPar
This is a manually\sphinxhyphen{}numbered footnote reference.%
\begin{footnote}[3]\sphinxAtStartFootnote
This is a manually\sphinxhyphen{}numbered footnote definition.
%
\end{footnote}

\item {} 
\sphinxAtStartPar
This is an auto\sphinxhyphen{}numbered footnote reference.%
\begin{footnote}[1]\sphinxAtStartFootnote
This is an auto\sphinxhyphen{}numbered footnote definition.
%
\end{footnote}

\end{itemize}


\bigskip\hrule\bigskip



\section{嵌入图像}
\label{\detokenize{other/myst_markdown:id18}}
\sphinxAtStartPar
要添加图像，请使用感叹号 (!), 然后在方括号增加替代文本，图片链接放在圆括号里，括号里的链接后可以增加一个可选的图片标题文本。

\begin{sphinxVerbatim}[commandchars=\\\{\}]
![\PYG{n+nt}{这是图片}](\PYG{n+na}{\PYGZus{}static/philly\PYGZhy{}magic\PYGZhy{}garden.jpg \PYGZdq{}Magic Gardens\PYGZdq{}})

    ```\PYGZob{}image\PYGZcb{} \PYGZus{}static/philly\PYGZhy{}magic\PYGZhy{}garden.jpg
    :alt: philly magic garden
    :class: bg\PYGZhy{}primary
    :scale: 50 \PYGZpc{}
    :align: center
    ```
\end{sphinxVerbatim}

\sphinxAtStartPar
显示效果

\sphinxAtStartPar
\sphinxincludegraphics{{philly-magic-garden}.jpg}

\noindent{\hspace*{\fill}\sphinxincludegraphics[scale=0.5]{{philly-magic-garden}.jpg}\hspace*{\fill}}


\bigskip\hrule\bigskip



\section{超链接}
\label{\detokenize{other/myst_markdown:id19}}
\sphinxAtStartPar
超链接Markdown语法代码：\sphinxcode{\sphinxupquote{{[}超链接显示名{]}(超链接地址 "超链接title")}}

\begin{sphinxVerbatim}[commandchars=\\\{\}]
这是一个链接 [\PYG{n+nt}{Markdown语法}](\PYG{n+na}{https://markdown.com.cn})
\end{sphinxVerbatim}

\sphinxAtStartPar
显示效果

\sphinxAtStartPar
这是一个链接 \sphinxhref{https://markdown.com.cn}{Markdown语法}


\bigskip\hrule\bigskip



\section{表格}
\label{\detokenize{other/myst_markdown:id20}}
\sphinxAtStartPar
表格必须使用MyST Markdown格式，否则PDF无法显示。

\sphinxAtStartPar
list\sphinxhyphen{}table 指令用于根据统一的两级项目符号列表中的数据创建表格。 “统一”是指每个子列表（二级列表）必须包含相同数量的列表项。

\begin{sphinxVerbatim}[commandchars=\\\{\}]
:::\PYGZob{}list\PYGZhy{}table\PYGZcb{} 表格标题
:widths: 15 10 30
:header\PYGZhy{}rows: 1

\PYG{k}{*}\PYG{+w}{ }  \PYGZhy{} \PYG{g+gs}{**Treat**}
\PYG{+w}{    }\PYG{k}{\PYGZhy{}}\PYG{+w}{ }\PYG{g+gs}{**Quantity**}
\PYG{+w}{    }\PYG{k}{\PYGZhy{}}\PYG{+w}{ }\PYG{g+gs}{**Description**}
\PYG{k}{*}\PYG{+w}{ }  \PYGZhy{} Albatross
\PYG{+w}{    }\PYG{k}{\PYGZhy{}}\PYG{+w}{ }2.99
\PYG{+w}{    }\PYG{k}{\PYGZhy{}}\PYG{+w}{ }On a stick!
\PYG{k}{*}\PYG{+w}{ }  \PYGZhy{} Crunchy Frog
\PYG{+w}{    }\PYG{k}{\PYGZhy{}}\PYG{+w}{ }1.49
\PYG{+w}{    }\PYG{k}{\PYGZhy{}}\PYG{+w}{ }If we took the bones out, it wouldn\PYGZsq{}t be
 crunchy, now would it?
\PYG{k}{*}\PYG{+w}{ }  \PYGZhy{} Gannet Ripple
\PYG{+w}{    }\PYG{k}{\PYGZhy{}}\PYG{+w}{ }1.99
\PYG{+w}{    }\PYG{k}{\PYGZhy{}}\PYG{+w}{ }On a stick!
:::
\end{sphinxVerbatim}

\sphinxAtStartPar
显示效果

\sphinxAtStartPar
:::\{list\sphinxhyphen{}table\} 表格标题
:widths: 15 10 30
:header\sphinxhyphen{}rows: 1
\begin{itemize}
\item {} \begin{itemize}
\item {} 
\sphinxAtStartPar
\sphinxstylestrong{Treat}

\item {} 
\sphinxAtStartPar
\sphinxstylestrong{Quantity}

\item {} 
\sphinxAtStartPar
\sphinxstylestrong{Description}

\end{itemize}

\item {} \begin{itemize}
\item {} 
\sphinxAtStartPar
Albatross

\item {} 
\sphinxAtStartPar
2.99

\item {} 
\sphinxAtStartPar
On a stick!

\end{itemize}

\item {} \begin{itemize}
\item {} 
\sphinxAtStartPar
Crunchy Frog

\item {} 
\sphinxAtStartPar
1.49

\item {} 
\sphinxAtStartPar
If we took the bones out, it wouldn’t be
crunchy, now would it?

\end{itemize}

\item {} \begin{itemize}
\item {} 
\sphinxAtStartPar
Gannet Ripple

\item {} 
\sphinxAtStartPar
1.99

\item {} 
\sphinxAtStartPar
On a stick!
:::

\end{itemize}

\end{itemize}


\bigskip\hrule\bigskip



\chapter{Indices and tables}
\label{\detokenize{index:indices-and-tables}}\begin{itemize}
\item {} 
\sphinxAtStartPar
\DUrole{xref,std,std-ref}{genindex}

\item {} 
\sphinxAtStartPar
\DUrole{xref,std,std-ref}{modindex}

\item {} 
\sphinxAtStartPar
\DUrole{xref,std,std-ref}{search}

\end{itemize}



\renewcommand{\indexname}{索引}
\printindex
\end{document}